\documentclass[aspectratio=169,10pt]{beamer}

\usepackage{fontspec}
\usepackage{microtype}
\usepackage{polyglossia}
\setmainlanguage{russian}
\setmainfont{Liberation Sans}[Ligatures=TeX]
\setsansfont{Liberation Sans}
\setmonofont{Liberation Mono}
\newfontfamily\cyrillicfont{Liberation Sans}
\newfontfamily\cyrillicfontsf{Liberation Sans}
\newfontfamily\cyrillicfonttt{Liberation Mono}

\usepackage{graphicx}
\usepackage{minted}
\usepackage{xcolor}
\usepackage{tikz}
\usetikzlibrary{arrows.meta,positioning}

\definecolor{Bg}{HTML}{2D353B}
\definecolor{Panel}{HTML}{343F44}
\definecolor{PanelAlt}{HTML}{3D484D}
\definecolor{Text}{HTML}{D3C6AA}
\definecolor{Muted}{HTML}{859289}
\definecolor{Green}{HTML}{A7C080}
\definecolor{Blue}{HTML}{7FBBB3}
\definecolor{Yellow}{HTML}{DBBC7F}
\definecolor{Red}{HTML}{E67E80}
\definecolor{Purple}{HTML}{D699B6}

\setbeamercolor{normal text}{fg=Text,bg=Bg}
\setbeamercolor{title}{fg=Green,bg=Bg}
\setbeamercolor{frametitle}{fg=Text,bg=Bg}
\setbeamercolor{structure}{fg=Blue}
\setbeamercolor{footline}{fg=Muted,bg=Bg}
\setbeamerfont{title}{size=\fontsize{30}{34}\selectfont,series=\bfseries}
\setbeamerfont{subtitle}{size=\fontsize{12}{16}\selectfont}
\setbeamerfont{frametitle}{size=\fontsize{18}{22}\selectfont,series=\bfseries}
\setbeamertemplate{navigation symbols}{}
\setbeamertemplate{itemize item}{\textcolor{Yellow}{\scriptsize$\blacksquare$}}
\setbeamertemplate{itemize subitem}{\textcolor{Blue}{\scriptsize$\bullet$}}
\setbeamertemplate{frametitle}{}
\setbeamertemplate{footline}{
  \hbox{
    \begin{beamercolorbox}[wd=.49\paperwidth,ht=2.5ex,dp=1.2ex,leftskip=0.35cm]{footline}
      FrogOS / reconciliation runtime
    \end{beamercolorbox}
    \begin{beamercolorbox}[wd=.49\paperwidth,ht=2.5ex,dp=1.2ex,rightskip=0.35cm plus1fil]{footline}
      \hfill \insertframenumber/\inserttotalframenumber
    \end{beamercolorbox}
  }
}

\setminted{
  bgcolor=Panel,
  breaklines=true,
  fontsize=\footnotesize,
  frame=none,
  linenos=false,
  tabsize=2
}
\usemintedstyle{gruvbox-dark}

\newcommand{\kicker}[1]{\textcolor{Muted}{\footnotesize\MakeUppercase{#1}}}
\newcommand{\slidehead}[2]{
  \vspace*{0.2cm}
  \kicker{#1}\par
  \vspace{0.1cm}
  {\usebeamerfont{frametitle}\textcolor{Text}{#2}}\par
  \vspace{0.18cm}
  {\color{Muted}\rule{\linewidth}{0.4pt}}
  \vspace{0.25cm}
}
\newcommand{\bigclaim}[1]{{\fontsize{22}{27}\selectfont\bfseries #1}}
\newcommand{\leadclaim}[1]{{\fontsize{18}{23}\selectfont\bfseries #1}}
\newcommand{\muted}[1]{\textcolor{Muted}{#1}}
\newcommand{\green}[1]{\textcolor{Green}{#1}}
\newcommand{\blue}[1]{\textcolor{Blue}{#1}}
\newcommand{\yellow}[1]{\textcolor{Yellow}{#1}}
\newcommand{\red}[1]{\textcolor{Red}{#1}}
\newcommand{\panel}[2]{
  \begin{beamercolorbox}[wd=#1,sep=0.25cm,rounded=false]{normal text}
    \colorbox{Panel}{\begin{minipage}{\dimexpr#1-0.5cm\relax}#2\end{minipage}}
  \end{beamercolorbox}
}
\newcommand{\pill}[2]{%
  \tikz[baseline=(x.base)] \node[inner xsep=6pt,inner ysep=2.5pt,fill=#1,text=Bg,font=\bfseries\scriptsize] (x) {#2};%
}

\tikzset{
  flow/.style={
    rectangle,
    rounded corners=2pt,
    minimum width=2.0cm,
    minimum height=0.9cm,
    align=center,
    draw=#1,
    fill=Panel,
    text=Text,
    line width=0.8pt,
    font=\small\bfseries
  },
  wire/.style={-{Latex[length=2.2mm]}, line width=0.7pt, draw=Muted}
}

\title{FrogOS}
\subtitle{declarative intent -> observed state -> ordered plan -> commit}
\author{Кроитор Александр}
\institute{Научный руководитель: Nichita Nartea}
\date{2026}

\begin{document}

\begin{frame}
  \vspace*{0.7cm}
  \kicker{курсовая работа / autonomous linux OS}
  \vspace{0.35cm}

  {\usebeamerfont{title}\textcolor{Green}{FrogOS}}\par
  \vspace{0.18cm}
  {\usebeamerfont{subtitle}\textcolor{Blue}{declarative intent -> observed state -> ordered plan -> commit}}\par

  \vfill
  \begin{tikzpicture}
    \node[flow=Green] (i) {Intent};
    \node[flow=Blue, right=0.7cm of i] (s) {Observed\\state};
    \node[flow=Yellow, right=0.7cm of s] (p) {Plan};
    \node[flow=Purple, right=0.7cm of p] (c) {Commit};
    \draw[wire] (i) -- (s);
    \draw[wire] (s) -- (p);
    \draw[wire] (p) -- (c);
  \end{tikzpicture}

  \vfill
  \muted{Кроитор Александр / научный руководитель: Nichita Nartea / 2026}
\end{frame}

\begin{frame}{}
  \slidehead{01 / problem}{Linux state drifts}

  \begin{columns}[T,totalwidth=\textwidth]
    \begin{column}{0.50\textwidth}
      \leadclaim{Админ описал \green{desired state}.}
      \vspace{0.45cm}

      \begin{itemize}
        \item package доступен;
        \item service включён;
        \item user состоит в группах;
        \item профиль реагирует на процесс.
      \end{itemize}
    \end{column}
    \begin{column}{0.44\textwidth}
      \leadclaim{Хост ушёл в \red{state drift}.}
      \vspace{0.45cm}

      \begin{itemize}
        \item пакет удалили;
        \item service упал;
        \item группа пропала;
        \item power profile сменился.
      \end{itemize}
    \end{column}
  \end{columns}

  \vfill
  {\large\yellow{Задача FrogOS: не выполнить скрипт, а вернуть систему к intent.}}
\end{frame}

\begin{frame}{}
  \slidehead{02 / requirements}{Функциональные требования}

  \begin{columns}[T,totalwidth=\textwidth]
    \begin{column}{0.48\textwidth}
      \green{Intent path}
      \begin{itemize}
        \item принять декларативную конфигурацию;
        \item проверить типы, семантику и бизнес логику;
        \item сохранить IR как \texttt{intent.json};
        \item явно версионировать контракт.
      \end{itemize}
    \end{column}
    \begin{column}{0.48\textwidth}
      \yellow{Runtime path}
      \begin{itemize}
        \item собрать \texttt{SystemState};
        \item построить delta между графами;
        \item применить ordered plan;
        \item отправить commit или rollback.
      \end{itemize}
    \end{column}
  \end{columns}

  \vfill
  \muted{Функционально FrogOS должна закрывать полный путь: describe -> observe -> plan -> execute -> store.}
\end{frame}

\begin{frame}{}
  \slidehead{03 / constraints}{Нефункциональные требования}

  \begin{columns}[T,totalwidth=\textwidth]
    \begin{column}{0.31\textwidth}
      \pill{Green}{Safety}\\[0.18cm]
      undo для каждого шага, dry-run путь, daemon boundary
    \end{column}
    \begin{column}{0.31\textwidth}
      \pill{Yellow}{Determinism}\\[0.18cm]
      одинаковые входы дают одинаковый план
    \end{column}
    \begin{column}{0.31\textwidth}
      \pill{Blue}{Audit}\\[0.18cm]
      поколение хранит проверенный intent и результат применения
    \end{column}
  \end{columns}

  \vspace{0.55cm}
  \begin{columns}[T,totalwidth=\textwidth]
    \begin{column}{0.48\textwidth}
      \green{Maintainability}
      \begin{itemize}
        \item новые домены проходят DSL -> IR -> Toad -> Planner -> Executor;
        \item side effects локализованы в backend traits.
      \end{itemize}
    \end{column}
    \begin{column}{0.48\textwidth}
      \yellow{Portability}
      \begin{itemize}
        \item сборка через flakes;
        \item e2e в FHS-контейнере;
        \item Rust/Haskell boundary через protobuf.
      \end{itemize}
    \end{column}
  \end{columns}
\end{frame}

\begin{frame}{}
  \slidehead{04 / model}{Reconciliation loop}

  \vspace{0.6cm}

  \begin{center}
    {\fontsize{24}{30}\selectfont
      \( I \times S^{obs}_t \rightarrow P_t \rightarrow S^{obs}_{t+1} \)
    }
  \end{center}

  \vspace{0.55cm}

  \begin{columns}[T,totalwidth=0.88\textwidth]
    \begin{column}{0.28\textwidth}
      \pill{Green}{I}\\[0.18cm]
      проверенный intent
    \end{column}

    \begin{column}{0.28\textwidth}
      \pill{Blue}{Sobs}\\[0.18cm]
      то, что увидел Toad
    \end{column}

    \begin{column}{0.28\textwidth}
      \pill{Yellow}{P}\\[0.18cm]
      ordered plan + undo
    \end{column}
  \end{columns}

  \vfill

  \muted{FrogOS не знает абсолютную истину о машине. Каждый цикл начинается с нового наблюдения.}
\end{frame}

\begin{frame}{}
  \slidehead{05 / architecture}{Компоненты не смешивают ответственность}

  \centering
  \begin{tikzpicture}[node distance=0.55cm]
    \node[flow=Green] (dsl) {DSL\\Haskell};
    \node[flow=Green, right=of dsl] (ir) {IR v9\\JSON};
    \node[flow=Blue, right=of ir] (toad) {Toad\\Rust probes};
    \node[flow=Yellow, right=of toad] (planner) {Planner\\Haskell};
    \node[flow=Purple, right=of planner] (exec) {Executor\\Rust};

    \draw[wire] (dsl) -- (ir);
    \draw[wire] (ir) -- (toad);
    \draw[wire] (toad) -- (planner);
    \draw[wire] (planner) -- (exec);
  \end{tikzpicture}

  \vspace{0.55cm}
  \begin{columns}[T,totalwidth=\textwidth]
    \begin{column}{0.47\textwidth}
      \green{Чистые слои}
      \begin{itemize}
        \item DSL validation;
        \item IR schema version;
        \item graph delta;
        \item JSON/protobuf contracts.
      \end{itemize}
    \end{column}
    \begin{column}{0.47\textwidth}
      \yellow{Side effects}
      \begin{itemize}
        \item Nix package store;
        \item dinit modules;
        \item users/groups;
        \item network rules.
      \end{itemize}
    \end{column}
  \end{columns}
\end{frame}

\begin{frame}[fragile]{}
  \slidehead{06 / DSL}{Intent читается как сценарий, но не является скриптом}

  \begin{minted}{haskell}
gaming :: ConfigBuilder ()
gaming = profile "gaming" $
    when (processRunning "steam" `via` poll 500) $ do
        disable nginx
        setPowerProfile Performance
  \end{minted}

  \vspace{0.25cm}
  \begin{columns}[T,totalwidth=\textwidth]
    \begin{column}{0.31\textwidth}
      \pill{Blue}{Condition}\\[0.15cm]
      когда профиль активен
    \end{column}
    \begin{column}{0.31\textwidth}
      \pill{Yellow}{Action}\\[0.15cm]
      какое состояние нужно
    \end{column}
    \begin{column}{0.31\textwidth}
      \pill{Purple}{Policy}\\[0.15cm]
      какие варианты допустимы
    \end{column}
  \end{columns}
\end{frame}

\begin{frame}[fragile]{}
  \slidehead{07 / comparison}{NixOS фиксирует систему, FrogOS реагирует на состояние}

  \begin{columns}[T,totalwidth=\textwidth]
    \begin{column}{0.48\textwidth}
      \green{NixOS}
      \begin{minted}[fontsize=\tiny]{nix}
{ pkgs, ... }: {
  services.nginx.enable = true;
  users.users.alice = {
    isNormalUser = true;
    extraGroups = [ "wheel" "docker" ];
  };
  environment.systemPackages = [
    pkgs.dust
  ];
}
      \end{minted}
      {\scriptsize\muted{Фичи: reproducible rebuild, generations, Nix store.}}
    \end{column}
    \begin{column}{0.48\textwidth}
      \yellow{FrogOS}
      \begin{minted}[fontsize=\tiny]{haskell}
profile "desktop" $ do
  enable nginx
  user "alice" $ do
    normalUser True
    extraGroups ["wheel", "docker"]
    exposePackage "dust"

profile "gaming" $
  when (processRunning "steam" `via` poll 500) $
    disable nginx
      \end{minted}
      {\scriptsize\muted{Фичи: runtime observation, reactive profiles, delta plan.}}
    \end{column}
  \end{columns}

  \vfill
  \muted{FrogOS не заменяет Nix. Он использует Nix как backend там, где нужны пакеты и store.}
\end{frame}

\begin{frame}[fragile]{}
  \slidehead{08 / IR}{JSON-контракт переживает runtime boundaries}

  \begin{minted}{json}
{
  "version": 9,
  "condition": {
    "type": "process_running",
    "name": "steam",
    "observe": { "type": "poll", "interval_ms": 500 }
  },
  "actions": [
    { "type": "module_disable", "priority": 100 },
    { "type": "power_profile", "value": "performance", "priority": 100 }
  ]
}
  \end{minted}

  \vspace{0.15cm}
  \muted{Версия схемы даёт Planner право явно отвергнуть несовместимый intent.}
\end{frame}

\begin{frame}{}
  \slidehead{09 / planner}{Planner считает delta, а не команды из профиля}

  \centering
  \begin{tikzpicture}[node distance=0.48cm]
    \node[flow=Green, minimum width=2.5cm] (active) {active\\profiles};
    \node[flow=Blue, right=of active, minimum width=2.5cm] (desired) {desired\\graph};
    \node[flow=Blue, below=0.55cm of desired, minimum width=2.5cm] (current) {current\\graph};
    \node[flow=Yellow, right=0.75cm of desired, minimum width=2.5cm] (delta) {delta};
    \node[flow=Purple, right=of delta, minimum width=2.5cm] (steps) {ordered\\PlanStep};

    \draw[wire] (active) -- (desired);
    \draw[wire] (current) -- (delta);
    \draw[wire] (desired) -- (delta);
    \draw[wire] (delta) -- (steps);
  \end{tikzpicture}

  \vspace{0.45cm}
  \begin{columns}[T,totalwidth=\textwidth]
    \begin{column}{0.48\textwidth}
      \green{Invariant}
      \begin{itemize}
        \item equal inputs -> equal plan;
        \item reconciled resource -> no step;
        \item conflict -> planning error.
      \end{itemize}
    \end{column}
    \begin{column}{0.48\textwidth}
      \yellow{Ordering}
      \begin{itemize}
        \item groups before users;
        \item users before memberships;
        \item packages before modules.
      \end{itemize}
    \end{column}
  \end{columns}
\end{frame}

\begin{frame}{}
  \slidehead{10 / planning logic}{Почему нужен граф ресурсов}

  \begin{columns}[T,totalwidth=\textwidth]
    \begin{column}{0.48\textwidth}
      \green{Desired graph}
      \begin{itemize}
        \item строится из IR и активных профилей;
        \item priority разрешает конкурирующие желания;
        \item effect resources не становятся состоянием сами по себе.
      \end{itemize}
    \end{column}
    \begin{column}{0.48\textwidth}
      \yellow{Current graph}
      \begin{itemize}
        \item строится из \texttt{SystemState};
        \item содержит только наблюдаемые факты;
        \item delta удаляет уже согласованные ресурсы.
      \end{itemize}
    \end{column}
  \end{columns}

  \vfill
  \begin{center}
    \pill{Blue}{Plan = minimal difference + backend-safe ordering + undo metadata}
  \end{center}
\end{frame}

\begin{frame}[fragile]{}
  \slidehead{11 / executor}{PlanStep несёт действие и обратный шаг}

  \begin{minted}{rust}
#[serde(tag = "type", rename_all = "snake_case")]
pub enum DeltaAction {
    GroupAdd { name: String },
    UserAdd { config: UserConfig },
    MembershipAdd { user: String, group: String },
    PackageInstall { name: String },
    ModuleEnable { module: ModuleRef, priority: u32 },
    NetworkApply { allow_ports: Vec<u16> },
}
  \end{minted}

  \vspace{0.2cm}
  \begin{center}
    \yellow{DeltaAction = данные. Side effect = только dispatch в Executor.}
  \end{center}
\end{frame}

\begin{frame}[fragile]{}
  \slidehead{12 / demo signal}{Реальный ordered plan заканчивается commit}

  \begin{minted}{text}
[1/11]  add group docker
[2/11]  add group networkmanager
[3/11]  add group nobody
[4/11]  add group nogroup
[5/11]  add group wheel
[6/11]  add user yorunikakeru
[7/11]  add yorunikakeru -> group docker
[8/11]  add yorunikakeru -> group networkmanager
[9/11]  add yorunikakeru -> group wheel
[10/11] install package dust
frogosd::supervisor_runtime: store commit
  \end{minted}

  \vspace{0.08cm}
  \muted{Лог показывает не “успешную команду”, а ordered reconciliation plan.}
\end{frame}

\begin{frame}{}
  \slidehead{13 / runtime}{Supervisor держит actor pipeline живым}

  \centering
  \begin{tikzpicture}[node distance=0.56cm]
    \node[flow=Blue, minimum width=2.35cm] (toad) {ToadActor};
    \node[flow=Green, right=of toad, minimum width=2.0cm] (bridge) {Bridge};
    \node[flow=Yellow, right=of bridge, minimum width=2.45cm] (planner) {PlannerActor};
    \node[flow=Purple, right=of planner, minimum width=2.45cm] (exec) {ExecutorActor};

    \draw[wire] (toad) -- (bridge);
    \draw[wire] (bridge) -- (planner);
    \draw[wire] (planner) -- (exec);

    \node[below=0.12cm of toad, text=Muted, font=\tiny] {publishes state};
    \node[below=0.12cm of bridge, text=Muted, font=\tiny] {encodes};
    \node[below=0.12cm of planner, text=Muted, font=\tiny] {plans};
    \node[below=0.12cm of exec, text=Muted, font=\tiny] {applies};

    \node[above=0.14cm of bridge, text=Blue, font=\tiny\bfseries] {watch};
    \node[above=0.14cm of planner, text=Blue, font=\tiny\bfseries] {mpsc + protobuf};
    \node[above=0.14cm of exec, text=Blue, font=\tiny\bfseries] {Plan};
  \end{tikzpicture}

  \vspace{0.28cm}
  \begin{columns}[T,totalwidth=\textwidth]
    \begin{column}{0.48\textwidth}
      \green{rest-for-one}
      \begin{itemize}
        \item Toad -> вся downstream chain;
        \item Bridge -> Planner + Executor;
        \item Executor -> новая execution chain.
      \end{itemize}
    \end{column}
    \begin{column}{0.48\textwidth}
      \yellow{Почему так}
      \begin{itemize}
        \item paired ends у каналов;
        \item state source должен быть цельным;
        \item restart rate limited.
      \end{itemize}
    \end{column}
  \end{columns}
\end{frame}

\begin{frame}{}
  \slidehead{14 / result}{Что защищает архитектура}

  \begin{columns}[T,totalwidth=\textwidth]
    \begin{column}{0.31\textwidth}
      \pill{Green}{Correctness}\\[0.2cm]
      типы DSL, IR v9, golden/round-trip tests
    \end{column}
    \begin{column}{0.31\textwidth}
      \pill{Yellow}{Recovery}\\[0.2cm]
      undo в каждом PlanStep, rollback в обратном порядке
    \end{column}
    \begin{column}{0.31\textwidth}
      \pill{Blue}{Isolation}\\[0.2cm]
      FHS e2e, daemon socket, generation store
    \end{column}
  \end{columns}

  \vspace{0.7cm}
  \begin{columns}[T,totalwidth=\textwidth]
    \begin{column}{0.55\textwidth}
      \leadclaim{\green{Итог:} self-healing как reconciliation loop.}
    \end{column}
    \begin{column}{0.39\textwidth}
      \begin{itemize}
        \item intent проверяется;
        \item state наблюдается;
        \item delta применяется с undo.
      \end{itemize}
    \end{column}
  \end{columns}
\end{frame}

\end{document}
